%====================================================================
\section{TeX Interface Macros (\texttt{tkz-elements.sty})}
\label{sec:sty_interface}

\subsection{Purpose}

The file \texttt{tkz-elements.sty} provides a small set of macros that
act as a bridge between \TeX{} and the Lua engine used by
\tkzNamePack{tkz-elements}. Their main goals are:

\begin{itemize}
\item printing Lua values safely in \TeX{};
\item transferring geometric objects (points, paths, circles) from Lua to TikZ;
\item simplifying the drawing stage when computations are done in Lua.
\end{itemize}

\medskip
Unless stated otherwise, these macros assume that \tkzNamePack{tkz-elements}
has been initialized in Lua (typically via |\directlua{init_elements()}|.

%--------------------------------------------------------------------
\subsection{Summary of provided macros}


\begin{tabular}{ll}
\toprule
\texttt{Macro} & \texttt{Reference} \\
\midrule
\tkzMacro{tkz-elements}{tkzUseLua\{...\}}               & \\
\tkzMacro{tkz-elements}{tkzPrintNumber\{...\}}          & \\
\tkzMacro{tkz-elements}{tkzEraseLuaObj\{name\}}         & \\
\tkzMacro{tkz-elements}{tkzDrawCoordinates(\dots)}      & \\
\tkzMacro{tkz-elements}{tkzGetPointsFromPath\{P\}\{A\}} & \\
\tkzMacro{tkz-elements}{tkzDrawPointsFromPath(\dots)}   & \\
\tkzMacro{tkz-elements}{tkzDrawSegmentsFromPaths(\dots)}& \\
\tkzMacro{tkz-elements}{tkzDrawCirclesFromPaths(\dots)} & \\
\bottomrule
\end{tabular}

\medskip
\texttt{Note:} The exact internal Lua object layout may evolve; these macros
are designed to remain stable at the user level.

%====================================================================
\subsection{Macro \tkzMacro{tkz-elements}{tkzUseLua}}
\label{subsec:tkzUseLua}

\medskip
\texttt{Syntax: }

\verb|tkzUseLua{<lua-expression>\}|


\medskip
\texttt{Description: }
Evaluates the Lua expression and prints its result in the \TeX{} stream.

\medskip
\texttt{Argument: }
\begin{itemize}
\item \tkzname{<lua-expression>}: a Lua expression returning a value.
\end{itemize}

\medskip
\texttt{Notes: }
\begin{itemize}
\item This macro is meant for \emph{values} (numbers, booleans, strings).
\item For Lua objects (point, line, circle, path, \dots), prefer dedicated
transfer/draw macros or call an explicit Lua method that returns a value.
\end{itemize}

\texttt{Example: }
\begin{minipage}{.52\textwidth}
\begin{verbatim}
\directlua{init_elements()
  z.A = point(0,0)
  z.B = point(4,2)
}
The slope is \tkzUseLua{(z.B.im-z.A.im)/(z.B.re-z.A.re)}.
\end{verbatim}
\end{minipage}

%%%%%%%%%%%%%%%%%%%%%%%%%%%%%%%

This macro evaluates a Lua expression and inserts its result in the document.


\texttt{Syntax: }\verb|\tkzUseLua{<lua-code>}|

\texttt{Example usage: }

\begin{mybox}
\begin{verbatim}
\tkzUseLua{checknumber(math.pi)}
\end{verbatim}
\end{mybox}


%====================================================================
\subsection{Macro \tkzMacro{tkz-elements}{tkzPrintNumber}}
This macro formats and prints a number using PGF’s fixed-point output, with a default precision of 2 decimal places.

\medskip

%\texttt{Syntax: }\verb|\tkzPrintNumber[<precision>]{<number>}|
\texttt{Syntax: }
\begin{quote}
\tkzMacro{tkz-elements}{tkzPrintNumber\{<lua-expression>\}} \qquad
\tkzMacro{tkz-elements}{tkzPN\{<lua-expression>\}}
\end{quote}

\texttt{Example: }

\begin{mybox}
\begin{verbatim}
\tkzPrintNumber{pi} % outputs 3.14
\tkzPrintNumber[4]{sqrt(2)} % outputs 1.4142
\end{verbatim}
\end{mybox}

\texttt{Alias: } The macro \tkzMacro{tkz-elements}{tkzPN} is a shorthand for \tkzMacro{tkz-elements}{tkzPrintNumber}.


\texttt{Example: }
\begin{minipage}{.52\textwidth}
\begin{verbatim}
\directlua{init_elements()
  x = math.pi/7
}
$\alpha = \tkzPN{x}$.
\end{verbatim}
\end{minipage}


%====================================================================
\subsection{Macro \tkzMacro{tkz-elements}{tkzEraseLuaObj}}

\label{sub:macro_tkzEraseLuaObj}

\texttt{Syntax: }
\begin{quote}
\tkzMacro{tkz-elements}{tkzEraseLuaObj\{<name>\}}
\end{quote}



\texttt{Description: }
Removes a stored Lua object from the global storage table (typically \code{z},
\code{L}, \code{C}, \code{CO}, etc., depending on your conventions).


\texttt{Syntax: } \verb|\tkzEraseLuaObj{<object>}|

\texttt{Example usage: }

\begin{mybox}
\begin{verbatim}
\tkzEraseLuaObj{z.A}
\tkzEraseLuaObj{T.ABC}
\end{verbatim}
\end{mybox}


\texttt{Notes: }
Use this macro to avoid name collisions when compiling large documents
or when reusing the same object names across figures.



%====================================================================
\subsection{Macro \tkzMacro{tkz-elements}{tkzPathCount}}
This macro retrieves the number of vertices in a Lua path and stores the result in a TeX macro.

\texttt{Syntax: }\verb|\tkzPathCount(<lua-path>){<MacroName>}|

\texttt{Arguments: }

\verb|<lua-path>|: identifier of a Lua \verb|path| object.
\verb|<MacroName>|: name of a TeX macro (without the backslash) that will store the count.


\texttt{Example usage: }

\begin{tkzexample}[latex=.5\textwidth]
\directlua{
 z.A = point(0, 0)
 z.B = point(5, 4)
 L.AB = line(z.A, z.B)
 PA.AB = L.AB:path(3)
 z.O = point(2, 0)
 }
\tkzPathCount(PA.AB){N}
\begin{center}
\begin{tikzpicture}
\tkzGetNodes
\tkzDrawCoordinates[blue](PA.AB)
\tkzDrawPointsFromPath[red,size=2](PA.AB)
  \foreach \i in {1,...,\N}{
 \tkzGetPointFromPath(PA.AB,\i){P\i}
 \tkzDrawSegment[orange](O,P\i)
  }
\end{tikzpicture}
\end{center}
\end{tkzexample}



%====================================================================
\subsection{Macro \tkzMacro{tkz-elements}{tkzDrawCoordinates}}
This macro draws a curve using the TikZ \code{plot coordinates} syntax, with coordinates generated by Lua.


\texttt{Syntax: } |\tkzDrawCoordinates[<options>](lua-path)|

\texttt{Example usage: }

%\begin{mybox}
\begin{tkzexample}[latex=.5\textwidth]
\directlua{
  z.A = point(0, 0)
  z.B = point(6, 0)
  z.C = point(1, 4)
  T.ABC = triangle(z.A, z.B, z.C)
  z.M = T.ABC.bc.mid
  z.N = T.ABC.ca.mid
  PA.path = T.ABC:path(z.M,z.N,4)
}
\begin{tikzpicture}[gridded]
\tkzGetNodes
\tkzDrawPolygon(A,B,C)
\tkzDrawCoordinates[orange,thick](PA.path)
\tkzDrawPoints(A,B,C,M,N)
\tkzLabelPoints(A,B)
\tkzLabelPoints[above](C)
\end{tikzpicture}
\end{tkzexample}
%\end{mybox}

\medskip
\texttt{Description: }
Draws a Lua \code{path} object as a TikZ path. This macro is typically used
after the computational phase in Lua. Draws the polyline defined by the sequence of points in a Lua path, using explicit \code{(x,y)} coordinates.


\medskip
\texttt{Alias: } % \let\tkzDrawPath\tkzDrawCoordinates
\label{subsec:tkzDrawPath}

\begin{quote}
\tkzMacro{tkz-elements}{tkzDrawPath[<tikz-options>](lua-path)}
\end{quote}

%====================================================================

\subsection{Macro \tkzMacro{tkz-elements}{tkzDrawPointsFromPath}}
\label{sub:macro_tkzDrawPointsFromPath}

This macro draws all points of a path using \tkzname{tkzDrawPoint}, without exporting point names.

\medskip
\texttt{Syntax: }
\verb|\tkzDrawPointsFromPath[<options>](lua-path)|

\medskip
\texttt{Example usage: }

\begin{tkzexample}[latex=.5\textwidth]
\directlua{
init_elements()
LP.test = list_point()
LP.test:add(point(0, 0))
LP.test:add(point(1, 0))
LP.test:add(point(1, 1))
PA.curve = LP.test:as_path()
}
\begin{tikzpicture}
\tkzDrawCoordinates[blue](PA.curve)
\tkzDrawPointsFromPath[red,size=2](PA.curve)
\end{tikzpicture}
\end{tkzexample}

%====================================================================


%====================================================================
\subsection{Macro \tkzMacro{tkz-elements}{tkzGetPointsFromPath}}
This macro defines TikZ points from a Lua path. Each point is named with a base name followed by an index.

\texttt{Syntax: }\verb|\tkzGetPointsFromPath[<options>](<path>,<basename>)|

\medskip
\texttt{Description: }
Exports the points of a Lua path as named \TeX/TikZ points.
The naming scheme is \code{<prefix>1}, \code{<prefix>2}, \dots

\medskip
\texttt{Example: }

\begin{tkzexample}[latex=.5\textwidth]
 \directlua{
  init_elements()
  PA.g = path({ "(0, 0)", "(1, 0)","(1, 1)" })
}
\begin{tikzpicture}
  \tkzGetPointsFromPath(PA.g,A)
  \tkzDrawPolygon(A_1,A_2,A_3)
  \tkzDrawPoints(A_1,A_2,A_3)
\end{tikzpicture}
\end{tkzexample}

%====================================================================




\subsection{Macro \tkzMacro{tkz-elements}{tkzGetPointFromPath}}
This macro extracts the i-th vertex of a Lua path and defines it as a TikZ point with a given name.

\texttt{Syntax: }\verb|\tkzGetPointFromPath(<PathLua>,<index>){<PointName>}|

\texttt{Arguments: }

\verb|<PathLua>|: identifier of a Lua \verb|path| object.
\verb|<index>|: index of the vertex to retrieve (TeX number or numeric macro).
\verb|<PointName>|: name of the TikZ point to create.
\medskip


\texttt{Example usage:}

\begin{mybox}
\begin{verbatim}
% Get the 3rd vertex of PA.A and name it P3
\tkzGetPointFromPath(PA.A,3){P3}
% With counting and a loop
\tkzPathCount(PA.A){N}
\foreach \i in {1,...,\N}{
\expandafter\tkzGetPointFromPath\expandafter(PA.A,\i){P\i}
}
\end{verbatim}
\end{mybox}
%====================================================================




\subsection{Macro \tkzMacro{tkz-elements}{tkzDrawSegmentsFromPaths}}
\label{sub:macro_tkzDrawSegmentsFromPaths}

\texttt{Description: }
Draws segments encoded in one or more Lua paths (pairs of points or consecutive
points, depending on the chosen convention).

\texttt{Syntax: }\verb|\tkzDrawSegmentsFromPaths[<options>](<lua-path-start>,<lua-path-end>)|

\texttt{Example: }

\begin{tkzexample}[vbox]
\directlua{
 init_elements()
 z.A = point(0, 0)
 z.B = point(3, 0)
 C.AB = circle(z.A, z.B)
 PA.c = path()
 PA.u = path()
 PA.v = path()
 for i = 0, 18, 2 do
   local angle = i * 2 * math.pi / 18
   local p = point(polar(8, angle))
   local Lu, Lv = C.AB:tangent_from(p)
   local u , v  = Lu.pb, Lv.pb
   PA.c:add_point(p)
   PA.u:add_point(u)
   PA.v:add_point(v)
 end}

\begin{center}
 \begin{tikzpicture}[scale =.6]
  \tkzGetNodes
  \tkzDrawSegmentsFromPaths[draw,red](PA.c,PA.u)
  \tkzDrawSegmentsFromPaths[draw,red](PA.c,PA.v)
  \tkzDrawCircle[blue](A,B)
  \tkzDrawPoints(A,I,u,v)
\end{tikzpicture}
 \end{center}
\end{tkzexample}
%====================================================================
\subsection{Macro \tkzMacro{tkz-elements}{tkzDrawCirclesFromPaths}}
\label{sub:macro_tkzDrawCirclesFromPaths}
This macro draws circles defined by two paths: one for the centers and one for the points through which each circle passes.

\medskip
\texttt{Syntax: }\verb|\tkzDrawCirclesFromPaths[<options>](<PA.center>,<PA.through>)|

\medskip
\texttt{Example usage: }

\begin{mybox}
\begin{verbatim}
\tkzDrawCirclesFromPaths[blue](PA.O, PA.A)
\end{verbatim}
\end{mybox}

See the document \tkzname{Euclidean Geometry} presented in \href{http://altermundus.fr}{altermundus.fr}for other examples.

\medskip
\begin{tkzexample}[vbox]
\directlua{
 init_elements()
 z.A = point(-4, -4)
 z.B = point(4, -4)
 L.AB = line(z.A, z.B)
 C.AB = circle(z.A, z.B)
 S.AB = L.AB:square()
 _, _, z.C, z.D = S.AB:get()
 z.O = S.AB.ac.mid
 z.a = S.AB.ab.mid
 z.c = S.AB.cd.mid
 z.o1 = tkz.midpoint(z.a, z.O)
 z.o = tkz.midpoint(z.c, z.O)
 z.b = S.AB.bc.mid
 z.g = (z.o1 - z.O) + (z.b - z.O)
 z.o2 = intersection(line(z.c, z.g), line(z.O,z.b))
 z.o3 = (z.o - z.O) + (z.o2 - z.O)
 z.j = intersection(line(z.c, z.b), line(z.o2,z.o3))
 local C2 = circle(z.O, z.a)
 local C1 = circle(z.o, z.c)
 local C3 = circle(z.o3, z.j)
 PA.pc, PA.pt, n = C1:CCC(C2, C3)
}

\begin{center}
\begin{tikzpicture}[scale=1]
\tkzGetNodes
\tkzDrawPolygon[cyan](A,B,C,D)
\tkzDrawCircles(O,a o,c o1,a o2,b o3,j)
\tkzDrawCirclesFromPaths[draw,red](PA.pc,PA.pt)
\end{tikzpicture}
\end{center}
\end{tkzexample}


%===============================================================

\subsection{Macro \tkzMacro{tkz-elements}{tkzDrawFromPointToPath}}
This macro draws a set of line segments from a given TikZ point to every point of a Lua \emph{path} object.

\medskip
\texttt{Syntax: }\verb|\tkzDrawFromPointToPath[<options>](<point-tikz>,<lua-path>)|

\medskip
\texttt{Arguments: }

\verb|<options>| (optional): pgf/tikz drawing options (color, thickness, etc.).

\verb|<point-tikz>|: the name of a TikZ point (already defined).

\verb|<lua-path>|: identifier of a Lua \verb|path| object.

\medskip
\texttt{Example usage: }

\medskip
\begin{tkzexample}[vbox]
\directlua{
 init_elements()
 z.O = point(0, 0)
 z.A = point(5, 1)
 z.Ta = point(8, 2)
 C.A = circle(z.A, z.Ta)
 L.OA = line(z.O,z.A)
 z.x, z.y = intersection(C.A, L.OA,{near=z.O})
 C.A.through = z.y
 z.m = tkz.midpoint(z.O, z.x)
 z.n = tkz.midpoint(z.O, z.y)
 z.w = tkz.midpoint(z.m, z.n)
 PA.A = path()
 PA.c = path()
 for t = 0, 1 - 1e-12, 0.1 do
   PA.A:add_point(C.A:point(t))
   local m = tkz.midpoint(z.O, C.A:point(t))
   PA.c:add_point(m)
 end
}
 \begin{center}
  \begin{tikzpicture}[scale=.8]
   \tkzGetNodes
   \tkzDrawCircle(A,Ta)
   \tkzDrawLine(O,A)
   \tkzDrawPoints(O,A,m,n,w,x,y)
   \tkzLabelPoints(O,A,m,n,w)
   \tkzDrawCircle[red](w,n)
   \tkzDrawPointsFromPath[red](PA.c)
   \tkzDrawPointsFromPath[blue](PA.A)
   \tkzDrawSegmentsFromPaths[draw,teal](PA.c,PA.A)
   \tkzDrawFromPointToPath[orange](O,PA.A)
 \end{tikzpicture}
  \end{center}
\end{tkzexample}



\subsection{Archimedes spiral, a complete example}

This example is built with the help of some of the macros presented in this section. It is inspired by an example created with tkz-euclide by Jean-Marc Desbonnez.

\vspace{1em}
\begin{tkzexample}[vbox]
\directlua{
  z.O = point(0, 0)
  z.A = point(5, 0)
  L.OA = line(z.O, z.A)
  PA.spiral = path()
  PA.center = path()
  PA.through = path()
  PA.ray = path()
  for i = 1, 48 do
      local k = i / 48
      z["P"..i] = L.OA:point(k)
      local angle = i * tkz.tau / 48
      z["r"..i] = point(polar(5, angle))
      ray = line(z.O, z["r"..i])
      circ = circle(z.O, z["P"..i])
      z["X"..i], _ = intersection(ray, circ, { near = z["r"..i] })
      PA.spiral:add_point(z["X"..i])
      PA.center:add_point(z.O)
      PA.through:add_point(z["P"..i])
      PA.ray:add_point(z["r"..i])
  end}
\begin{center}
  \begin{tikzpicture}[scale = .8]
  \tkzGetNodes
  \tkzDrawCircle(O,A)
  \tkzDrawCirclesFromPaths[draw,
      orange!50](PA.center,PA.through)
  \tkzDrawSegmentsFromPaths[draw,teal](PA.center,PA.ray)
  \tkzDrawPath[red,thick,smooth](PA.spiral)
  \tkzDrawPointsFromPath[blue, size=2](PA.spiral)
  \end{tikzpicture}
\end{center}

\end{tkzexample}



\endinput

